\label{chap:taxonomia-monodromia}

La représentation positionnelle d'un nombre enregistre une suite de symboles,
mais ne conserve pas nécessairement l'état qui détermine sa continuation.
L'holographie modulaire triadique distingue les deux niveaux. La phase nonadique indique
la position d'une lecture dans le cycle de neuf pas ; l'état
enrichi conserve, en outre, l'orientation, le quotient euclidien, le transport
d'entraînement et la condition de frontière. Le retour de la phase n'implique donc pas
le retour de l'état.

Soit \(S\) un espace d'états enrichis, soit
\(T:S\to S\) une transition et soit
\[
g:S\longrightarrow\mathbb Z/9\mathbb Z
\]
la phase, avec \(g(Ts)=g(s)+1\). Pour une base \(b\geq2\), un lecteur
\(d:S\to\{0,\ldots,b-1\}\) produit
\[
x(s)=\sum_{n\geq0}\frac{d(T^ns)}{b^{n+1}}.
\]
Lorsqu'une valeur admet les deux développements positionnels usuels, on adopte la
représentation canonique qui ne se termine pas par une queue constante \(b-1\). Cette
convention sépare une ambiguïté d'écriture de la multiplicité proprement
HMT, qui procède de sections enrichies distinctes ayant une même
évaluation.
L'opérateur de retour nonadique est la restriction
\[
M=T^9:g^{-1}(r)\longrightarrow g^{-1}(r).
\]
La différence entre \(T^9s=s\) et \(g(T^9s)=g(s)\) formalise la différence
entre périodicité et monodromie.

La transition de l'état enrichi de TPK peut être donnée par une relation et non par une
fonction. Dans ce cas, la notation précédente s'applique au décalage sur
l'espace des histoires infinies ou à une réalisation déterministe dont l'état
a été étendu avec l'information nécessaire pour fixer la continuation ; on ne
présuppose pas qu'un bloc visible possède un successeur unique.

\begin{definition}[Nombre HMT et présentation arithmétique]
Un nombre HMT est uniquement une section compatible
\(\boldsymbol x\) du système inverse d'états survivants. Un
lecteur positionnel \(L\) dont le domaine contient \(\boldsymbol x\) étant choisi, le couple
\((\boldsymbol x,L)\) est une présentation ou une mesure du nombre. La valeur
archimédienne publiée par cette présentation est l'intersection ponctuelle des
cylindres associés lorsque leurs diamètres tendent vers zéro. Deux sections ayant la
même valeur ne sont pas identifiées avant l'application du quotient archimédien.
\end{definition}

\begin{definition}[Périodicités visible et enrichie]
La sortie est éventuellement périodique lorsqu'il existe \(N\geq0\) et \(p>0\)
tels que
\[
d(T^{n+p}s)=d(T^ns),\qquad n\geq N.
\]
L'état est éventuellement périodique lorsque
\[
T^{N+p}s=T^Ns.
\]
La seconde condition implique la première. La réciproque exige que le lecteur
sépare les queues de l'orbite considérée, c'est-à-dire que
\[
 \bigl(d(T^{n+k}s)\bigr)_{k\geq0}
 =\bigl(d(T^{m+k}s)\bigr)_{k\geq0}
 \quad\Longrightarrow\quad T^ns=T^ms.
\]
\end{definition}

\begin{theorem}[Classification par retour de phase]
\label{thm:clasificacion-retorno-fase}
Soit \(s\in S\).
\begin{enumerate}
  \item Si l'orbite de \(s\) atteint un état dont la sortie future est
  constamment nulle, alors \(x(s)\) possède un développement fini et est
  rationnel.
  \item Si l'état est éventuellement périodique, alors la sortie est
  éventuellement périodique et \(x(s)\in\mathbb Q\).
  \item Si le lecteur sépare les queues de l'orbite et que l'orbite de \(s\) n'est pas
  éventuellement périodique, alors la sortie n'est pas éventuellement périodique
  et \(x(s)\notin\mathbb Q\).
  \item La non-périodicité ne distingue pas l’irrationalité algébrique de la
transcendance. Cette distinction exige un certificat polynomial ou un
théorème de transcendance supplémentaire.
\end{enumerate}
\end{theorem}

\begin{proof}
Les deux premiers points sont la caractérisation classique des rationnels
par des développements positionnels finis ou ultimement périodiques. Dans
le troisième, une périodicité ultime de la sortie produirait deux queues
égales. L’hypothèse de séparation identifierait alors les états
correspondants, ce qui contredirait la non-périodicité de l’orbite. Le quatrième
point découle de l’existence d’irrationnels algébriques et transcendants
dont les développements ne sont pas ultimement périodiques.
\end{proof}

\begin{corollary}[Monodromie nonadique et apériodicité]
Supposons que la phase revienne tous les neuf pas, que l’orbite de retour
\[
 \{M^ks:k\geq0\},\qquad M=T^9,
\]
ne soit pas ultimement périodique et que le lecteur sépare les queues. Alors la
sortie n’est pas ultimement périodique et le nombre produit est irrationnel. La
périodicité de la phase coexiste dans ce cas avec une évolution apériodique de
l’état enrichi.
\end{corollary}

L’hypothèse ne peut être remplacée par une seule inégalité \(M(s)\ne s\), ni
par un nombre fini de retours distincts. Relativement à l’état initial,
à la transition et au lecteur déclarés, la fenêtre actuelle certifie neuf
retours de phase avec changement d’état ; elle démontre une monodromie finie, mais non
à elle seule l’apériodicité de toute l’orbite.

\ifdefined\HMTRepetirResultadosTaxonomicos
\section{Le carré local de \texorpdfstring{\(e\)}{e} et la rupture de mémoire}

Pour un mot décimal \(a_1\cdots a_n\), nous appellerons \emph{carré local}
de longueur \(\ell\) à la position \(i\) une factorisation
\[
 a_i\cdots a_{i+2\ell-1}=vv,
 \qquad |v|=\ell.
\]
Cette notion appartient à la combinatoire des mots et ne doit pas être confondue avec
une période de la queue infinie.

\begin{theorem}[Caractérisation finie du carré \(1828^2\)]
Considérons les \(243\) mots du catalogue N33 et le lecteur figé
\(w_6\mapsto w_{30}\mapsto d_1d_2d_3d_4\), avec quatre triades en base mille.
Il existe un unique mot dont la lecture contient un carré de longueur quatre
commençant au deuxième chiffre décimal. C’est
\[
 w_e=201101,
 \qquad
 718281828459=7\,\underbrace{1828\,1828}_{(1828)^2}\,459.
\]
La répétition ne se prolonge pas en une troisième copie : le chiffre suivant requis
serait \(1\), tandis que le chiffre effectif est \(4\). Dans tout le catalogue,
il n’existe qu’un autre carré de longueur quatre,
\[
 575157518472=\underbrace{5751\,5751}_{\text{carré de mot}}\,8472,
\]
mais il commence au premier chiffre.
\end{theorem}

\begin{proof}
L’énumération exhaustive évalue les \(243\) mots par arithmétique
entière exacte. Le profil des carrés pour les longueurs \(1,\ldots,6\) contient,
respectivement,
\[
240, 21, 3, 2, 0, 0
\]
occurrences, réparties entre
\[
153, 19, 3, 2, 0, 0
\]
mots. Les deux occurrences de longueur quatre sont exactement celles écrites
dans l’énoncé. Deux énumérations entières indépendantes reproduisent le même
recensement et leurs sorties complètes coïncident terme à terme.
\end{proof}

Le mot \(w_e\) possède en outre la chaîne de relèvement
\[
201101\mid121221\mid102011\mid012222\mid102011.
\]
Le troisième et le cinquième blocs visibles coïncident. En revanche,
les états qui les transportent ne coïncident pas. Avant les deux émissions, leurs réductions
modulo \(27\) sont
\[
(25,11,7,17,8,16),
\qquad
(7,8,4,23,26,7),
\]
et les quotients de Hensel ultérieurs sont
\[
(6,13,9,2,13,1),
\qquad
(15,0,3,19,23,25).
\]
Par conséquent, le retour \(102011\) est un retour de la projection visible, non
de l’état enrichi. C’est la formulation exacte de l’intuition d’une
répétition locale qui se rompt : le quotient mémorise une généalogie que le
bloc résiduel ne conserve pas.

Le résultat appartient au lecteur enrichi déjà fixé par le système
APP--TRIT--TPK\@. Le recensement fini caractérise exhaustivement l’apparition du
carré local ; le prolongement coinductif de la section appartient à la
connexion nonadique construite dans le
\cref{chap:numero-heisenberg-d108}. La propriété combinatoire du préfixe et
la génération du développement complet sont donc deux niveaux
compatibles d’une même généalogie.

Cette formulation fournit le cadre exact pour interpréter le début
décimal de \(e\),
\[
e=2.718281828459\ldots.
\]
L’expression correcte n’est plus « semi-périodicité », mais \emph{carré local
de mot avec retour visible et rupture de mémoire}.

\section{2 mesures finies de l’émetteur \(468\to243\) : fréquence et priorité structurelle}

Il n’existe pas de distribution uniforme sur tous les nombres réels. Il
existe en revanche une mesure canonique sur le catalogue fini lorsque l’on munit d’une loi uniforme
l’ensemble des \(104\,976\) germes TPK\@. Pour un mot visible \(w\), on
définit
\[
\mu_{\mathrm{TPK}}(w)
=\frac{\#\{\text{germes se projetant sur }w\}}{104\,976}.
\]
La mesure n’est pas uniforme. Les \(243\) mots possèdent des poids compris
entre \(144\) et \(2088\), avec neuf valeurs de multiplicité distinctes.
L’histogramme exact, où le poids est écrit comme base et le nombre de mots
comme exposant, est
\[
144^{120},\ 288^{54},\ 432^{18},\ 576^9,\ 864^{15},\
1008^6,\ 1440^3,\ 1944^{12},\ 2088^6.
\]

Pour les trois mots distingués, on obtient
\[
\mu_{\mathrm{TPK}}(010211)=\frac{1008}{104\,976}=\frac7{729},
\]
\[
\mu_{\mathrm{TPK}}(201101)
=\mu_{\mathrm{TPK}}(121200)
=\frac{144}{104\,976}=\frac1{729}.
\]
Le mot associé à \(\pi\) a cinq antécédents \(U_6\) et un poids sept
fois supérieur à celui des mots associés à \(e\) et \(\varphi\). Ce résultat
démontre une hiérarchie combinatoire dans le catalogue ; il n’identifie pas à lui
seul les mots aux constantes classiques et n’établit pas de probabilité
sur \(\mathbb R\).

Il ne définit pas non plus, à lui seul, une priorité totale : six mots ont un poids
\(1008\), tandis que \(e\) et \(\varphi\) appartiennent à la classe minimale de
poids \(144\), partagée par \(118\) autres mots. La distinction HMT doit
s’exprimer par un profil structurel, par exemple
\[
\mathcal P(w)=
\bigl(m_{\rm semilla}(w),m_{U_6}(w),r_{\rm supervivencia}(w),
s_{\rm frontera}(w),s_{\rm excepcional}(w)\bigr),
\]
où chaque coordonnée est définie avant l’application d’étiquettes externes. Un scalaire
de « priorité » exige en outre de fixer, également à l’avance, la manière d’ordonner ou de
pondérer ces coordonnées.
\fi

\section{Classification des constantes selon l’antécédent généalogique, l’opération génératrice et le codomaine}

La classification HMT sépare des propriétés que la représentation ordinaire
réunit sous une seule valeur :
\begin{center}
\begin{tabular}{p{0.25\textwidth}p{0.64\textwidth}}
\toprule
Classe & Condition structurelle \\
\midrule
Fini & La sortie atteint une queue nulle. \\
Rationnel périodique & La sortie est ultimement périodique. \\
Rationnel monodromique & La sortie est périodique, mais l’état enrichi
parcourt une fibre non triviale que le lecteur visible oublie. \\
Irracional coinductivo & La sortie n’est pas ultimement périodique et les
cylindres déterminent une unique valeur. \\
Recurrente local & Un mot fini contient des retours ou des carrés, sans que
cela impose une périodicité à la queue. \\
Algébrique & La valeur satisfait un polynôme non nul à coefficients entiers ;
la route doit être accompagnée de ce certificat. \\
Transcendant & La valeur ne satisfait aucun polynôme entier non nul ; la
classification exige un théorème de transcendance. \\
Distinguido HMT & La section satisfait un prédicat structurel figé de
fibre, de frontière, de survie ou de symétrie ; c’est une classe transversale aux
précédentes. \\
\bottomrule
\end{tabular}
\end{center}

Ainsi, HMT ne modifie pas les définitions arithmétiques correctes de la rationalité,
de l’algébricité et de la transcendance. Elle les affine en ajoutant la généalogie de l’état,
la monodromie de phase et la multiplicité des fibres que l’évaluation
archimédienne avait oubliées.

En particulier, \(\varphi\) est algébrique parce qu’il satisfait
\(X^2-X-1=0\), tandis que la transcendance de \(e\) et de \(\pi\) découle de
théorèmes classiques externes. Ces propriétés sont compatibles avec le fait que les trois
sections soient distinguées par HMT\@. La clôture dodécaphasique construit plus
loin la projection rationnelle exacte
\(\alpha_{12}=\pi_{12}(\alpha_{\mathrm{HMT}})\), et le prolongement compatible
construit la section coinductive \(\alpha_{\mathrm{HMT}}\), dans le
\cref{chap:alpha-rev7} ; aucune d’elles n’est utilisée dans cette taxonomie comme entrée
ni publiée ici comme une seconde résidence. Le noyau \(E_{21/22}\)
produit une racine analytique distincte, dont la concordance avec la voie entière est
démontrée dans le codomaine de \(\operatorname{Tr}_9\). Section, projection
rationnelle, racine analytique, lecteur et valeur conservent ainsi leurs types propres.

Enfin, une « direction infinitésimale » exige un typage propre. Pour un
lecteur fini \(r_m\), deux sections distinctes dans une même fibre de \(r_m\)
définissent une différence invisible à la profondeur \(m\), quoique active dans un
prolongement ultérieur. Elles ne constituent pas pour autant un réel infinitésimal non nul.
Si deux sections coïncident sous toutes les troncatures fidèles du système
inverse, elles sont la même section. Un infinitésimal arithmétique véritable exigerait de
construire une extension ordonnée ou valuée non archimédienne et de démontrer que les
opérations HMT y descendent.
